% !TeX root = ../../surface-lagrangian-torus-fibration.tex
% 2026-11-21 / v1.1: source-free slab の体積を用いた中央全点の一様有界径．
\subsection{固定点を含む中央 source の全点接続}

\begin{lemma}[中央 source-free slab の体積と接続]
\label{ke:source-central-volume-anchor}
命題\ref{ke:source-shell-gaussian-tail}の source-free shell $\Sigma_t$ に対し，
その固定 chart 内に compact な集合 $G_t$ と $t$ に依らない $c,C>0$ を取って
\[
 \operatorname{Vol}_{g_t}(G_t)\ge cL_t^{-1},\qquad
 \sup_{q\in G_t}d_{g_t}(q,\Sigma_t)\le C,
 \qquad \operatorname{diam}_{g_t}\Sigma_t\le C
\]
とできる．
\end{lemma}

\begin{proof}
$\Sigma_t$ の residual 半径を $\beta$ とし，
$0<\beta_-<\beta_+<\beta$ を固定する．正しい $x$-chart で
\[
 G_t=\{\beta_-\le|z|\le\beta_+,\quad
                  L_t/4\le s_x=-\log|x|^2\le L_t/3,\quad y=tz/x\}
\]
と置き，全ての angles を含める．$s_y=L_t-s_x-\log|z|^2$ なので両側の normal 深度は
$L_t$ と比較可能であり，$G_t$ は小さい $t$ で固定 physical source chart 内に入る．
補題\ref{ke:source-tube-potential}から，この集合上で
\[
 \omega_t^2\asymp L_t^{-2}|d\log x\wedge dz|^2.
\]
一周期の $\arg x$ と幅 $L_t/12$ の $s_x$ 区間，正の固定面積を持つ residual annulus を
積分することで，$\operatorname{Vol}_{g_t}(G_t)\asymp L_t^{-1}$ を得る．
体積の固定規格化係数は定数に吸収する．

$G_t$ と $\Sigma_t$ を含む一つの固定 source-free residual annulus を取る．
束 frame の深度との差は有界で，normal annulus の長さは $M_t=L_t+O(1)$ である．
両集合では $\min(s,M_t-s)\ge cL_t$ だから，canonical moment の差商上界により
$|\mu_t|\le C/L_t$ となる．系\ref{ke:fixed-base-central-band}を，例えば
$m_t=L_t^{1/2}$ としてこの annulus に適用すると，
$G_t\cup\Sigma_t$ の actual ambient 径は $t$ に一様に有界である．
ここでは source を含む中央領域の径を入力にしていない．
\end{proof}

\begin{theorem}[中央 source の一様有界な ambient 径]
\label{ke:source-central-diameter}
本節の固定 source chart と固定半径 $c_1$，および
命題\ref{ke:source-shell-gaussian-tail}の source-free shell $\Sigma_t$ に対し，
十分小さい $t$ で
\begin{equation}\label{ke:eq:source-central-diameter}
 \sup_{p\in\overline{B_t(c_1)}}d_{g_t}(p,\Sigma_t)
 \le C,\qquad
 \operatorname{diam}_{d_{g_t}}\overline{B_t(c_1)}
 \le C
\end{equation}
が成立する．これは source 固定点と二本の座標軸を含む全点の評価であり，
正規化した ambient 径は $O(R_t^{-1})$ で消失する．
追加の tensor 比較仮定や potential の一様連続度は要求しない．
\end{theorem}

\begin{proof}
以下 $L=L_t$，$R=\log L$ と書き，actual Riemannian 体積を $\operatorname{Vol}_t$
で表す．K\"ahler 体積との固定係数は定数に吸収する．
補題\ref{ke:source-central-volume-anchor}により，
ある固定 $A_0,C_0,c_0>0$ と compact な集合 $G_t$ を選んで
\[
 G_t\subset\{d_{g_t}(\cdot,\Sigma_t)\le A_0\},\qquad
 \operatorname{Vol}_t(G_t)\ge c_0L^{-1},\qquad
 \operatorname{diam}_{d_{g_t}}\Sigma_t\le C_0
\]
とできる．$G_t$ は固定 source-free residual annulus 上の central normal slab である．
比較に用いる体積は $L^{-1}$ の尺度であり，以下の広い中央帯の Gaussian tail と一致する．

$p\in\overline{B_t(c_1)}$ とし，$D=d_{g_t}(p,\Sigma_t)$ と置く．
命題\ref{ke:source-potential-buffer}により，固定した $\delta,\kappa,b>0$ が存在して
\[
 B_{g_t}(p,\delta)\subset
 \mathcal D_{t,\kappa,b}
 :=\{s_x,s_y\ge\kappa L,\ |z|\le b\}
\]
となる．ここで $s_x=+\infty$ on $\{x=0\}$，$s_y=+\infty$ on $\{y=0\}$ とする．
この包含は actual ball 全体について成立し，rescaled bidisc 内への包含を
要求していない．この球の全ての点 $q$ は
$d_{g_t}(q,\Sigma_t)\ge(D-\delta)_+$ を満たすので，
命題\ref{ke:source-wide-gaussian-tail}から
\begin{equation}\label{ke:eq:source-ball-gaussian-upper}
 \operatorname{Vol}_t B_{g_t}(p,\delta)
 \le C L^{-1}\exp\bigl(-a(D-\delta)_+^2\bigr).
\end{equation}

一方，$\Sigma_t$ の径上界から
$G_t\subset B_{g_t}(p,D+A_0+C_0+1)$ である．
Compact な実四次元多様体 $(X_t,g_t)$ は complete であり，
$\operatorname{Ric}(g_t)=-g_t$ であるから，Bishop--Gromov 体積比較を
中心 $p$ の二つの球に適用できる．比較双曲空間の半径 $r\ge1$ の球体積は
$C\exp(Cr)$ 以下なので，定数を取り直して
\begin{equation}\label{ke:eq:source-ball-bishop-lower}
 \operatorname{Vol}_t B_{g_t}(p,\delta)
 \ge c L^{-1}\exp\bigl(-C(D+1)\bigr)
\end{equation}
を得る．大きい方の球や $G_t$ への接続が source chart 内に留まる必要はない．
全曲面の Ricci 下界と ambient 距離を同じ規約で用いている．

\eqref{ke:eq:source-ball-gaussian-upper}と
\eqref{ke:eq:source-ball-bishop-lower}の対数を取ると
\[
 a(D-\delta)_+^2\le C(D+1)+C.
\]
従って $D\le C$ であり，第一式を得る．
二点をそれぞれ $\Sigma_t$ に接続し，その ambient 径 $C_0$ を加えると
第二式も従う．全ての定数は $p,t$ に依存せず，測度零の点を除いていない．
\end{proof}

射影を保つ source 正規形では元の $x,y$ と新座標の差が $O(t)$ である．
従って固定 rescaled cores は半径を一度大きくして互いに含まれ，
上の全点径評価は命題\ref{ke:projected-source-normal-form}の座標にも適用できる．

\begin{corollary}[Source を含む中央 moment 帯の全点収縮]
\label{ke:source-central-band-collapse}
命題\ref{ke:source-central-band-reduction}の固定 source tube において，
$A>0$ と十分小さい固定残余半径 $b_0<c_1^2$ を固定する．
\[
 \mathcal E_{t,A,b_0}
 =\{|z|\le b_0,\ |\mu_t|\le A/L_t\}
\]
の全点は $\Sigma_t$ へ actual 長さ $C_A$ 以下で接続でき，
同集合の ambient 径も $C_A$ 以下である．
命題\ref{ke:fixed-stage-source-neck}の共通 normal Moser map の像についても，
両端の trajectory の長さを加えた同じ評価が成立する．
従ってこの中央帯に含まれる pinched fiber と，同じ範囲の複数 fibers の合併は
正規化した ambient 径で収縮する．
\end{corollary}

\begin{proof}
命題\ref{ke:source-central-band-reduction}は各点を $z$ を保って
$\overline{B_t(c_1)}$ へ $O_A(1)$ の費用で結ぶ．
前定理を続けて適用すればよい．Moser 像では
補題\ref{ke:short-circle-actual-buffer}の trajectory 評価を両端に加える．
これは ambient 径の主張であり，経路が一つの fiber 又は moment level の
中に留まることを要求しない．
\end{proof}

上の系は中央 moment 幅 $O(L_t^{-1})$ の範囲を扱う．
共通 interval 全体の同一 moment 退出は
定理\ref{ke:source-all-moment-exit}で別に証明する．
木上の fibers をどの target 座標へ対応させるか，および全曲面の no-shortcut 下界は
別の問題である．

\begin{corollary}[固定残余基底の source-inclusive 中央帯]
\label{ke:fixed-base-source-central-band}
命題\ref{ke:fixed-stage-source-neck}の一つの固定段階を取り，残余基底 $P$ を
連結 compact 曲面とする．$P$ 内の source は有限個とし，その周囲の互いに素な
固定小円板を除いた source-free 部分も連結に取る．
この neck の共通 canonical moment を $\mu_t$ とすると，任意の固定 $A>0$ に対し
\[
 \operatorname{diam}_{d_{g_t}}
 \{p\text{ in this neck}:|\mu_t(p)|\le A/L_t\}
 \le C_{P,A}
\]
となる．Actual normal Moser 像についても同じ結論が成立する．
従ってこの固定段階で構成した $\mu_t=0$ の閉じた pinched torus は，
source 部分だけでなく全体が正規化 ambient 径で収縮する．
\end{corollary}

\begin{proof}
各小円板では前の系を使い，その外側の固定 source-free shell へ接続する．
Shell 上の canonical moment は $O(L_t^{-1})$ である．
ここで source 座標と共通 neck の primitive の差は残余基底だけの項なので，
系\ref{ke:source-central-moment}により moment の加法定数は変わらない．
残る固定 source-free 基底と全ての shells を含む連結 compact 部分には，
定数 $A$ を一度大きくして系\ref{ke:fixed-base-central-band}を用いる．
Source 数は有限であり，各点からこの共通部分への費用は
$O_{P,A}(1)$，共通部分内の接続費用は $O_{P,A}(1)$ である．
Moser 像へは両端の短い trajectories を加える．
\end{proof}
